%!TEX root = ../main.tex
\section{Structural Indicators as Option Greeks}\label{sec:greeks}

The structural relation $E_t = e^{\pi_t}(A_t - D_t)^+$ endows MSTR equity with an
option-like payoff on Bitcoin.  If the NAV premium $\pi_t$ encodes the value of the
``accumulation option''---the right to issue equity at a premium and buy BTC
accretively---then the five indicators derived in Section~\ref{sec:indicators}
are precisely the option Greeks of this embedded claim.

This section formalises the mapping, derives each Greek from the structural model,
and introduces two new sensitivity measures.

%% ------------------------------------------------------------------
\subsection{The Accumulation Option}

\begin{definition}[Accumulation Option Value]\label{def:acc-option}
Define the \emph{accumulation option value} as the excess of equity market value
over net asset value:
$$
  V_{\mathrm{acc}}(t) \;=\; E_t - \mathrm{NAV}_t
  \;=\; \bigl(e^{\pi_t} - 1\bigr)(A_t - D_t)^+.
$$
When $\pi_t > 0$ the option is ``in the money'': the market assigns positive
value to management's ability to grow BTC holdings accretively.
When $\pi_t \le 0$ the option is worthless and equity trades at or below NAV.
\end{definition}

%% ------------------------------------------------------------------
\subsection{ILE as Delta}\label{ssec:ile-delta}

\begin{proposition}[ILE $\equiv$ Leverage Delta]\label{prop:ile-delta}
Let $E_t = e^{\pi_t}(H_t S_t - D_t)^+$ with $\pi_t$ held fixed.  Then
$$
  \Delta_{\mathrm{lev}}(t)
  \;=\;
  \frac{\partial \log E_t}{\partial \log S_t}\bigg|_{\pi_t}
  \;=\;
  \frac{A_t}{A_t - D_t}
  \;=\;
  \mathrm{ILE}_t.
$$
\end{proposition}

\begin{proof}
With $\pi_t$ fixed, $\log E_t = \pi_t + \log(H_t S_t - D_t)$.  Differentiating
with respect to $\log S_t$ (holding $H_t, D_t$ constant):
$$
  \frac{\partial \log E_t}{\partial \log S_t}
  = \frac{H_t S_t}{H_t S_t - D_t}
  = \frac{A_t}{A_t - D_t}
  = \beta_{LS}(t)
  = \mathrm{ILE}_t.
$$
\end{proof}

\paragraph{Economic intuition.}
Delta measures how much equity moves per unit BTC move.  Leverage amplifies
exposure exactly as a deep in-the-money call does: the higher the
debt-to-asset ratio, the larger $\Delta_{\mathrm{lev}}$.  With current
calibration,
$$
  \mathrm{ILE}_0 = 1.49,
$$
implying that a 1\% move in BTC produces approximately a 1.49\% move in
MSTR equity from balance-sheet leverage alone---comparable to a call with
$\Delta > 1$ when expressed in log-return terms.

%% ------------------------------------------------------------------
\subsection{TEE as Adjusted Delta}\label{ssec:tee-adj-delta}

\begin{proposition}[TEE $\equiv$ Total (Adjusted) Delta]\label{prop:tee-delta}
Let $\gamma_{\pi S}$ denote the empirical sensitivity of the log-premium to
BTC log-returns.  The total equity elasticity
$$
  \Delta_{\mathrm{tot}}(t)
  = \frac{d\log E_t}{d\log S_t}
  = \gamma_{\pi S} + \mathrm{ILE}_t
  = \mathrm{TEE}_t
$$
is the \emph{adjusted delta} that incorporates the premium channel.
\end{proposition}

\begin{proof}
From $\log E_t = \pi_t + \log(A_t - D_t)^+$, the total derivative with respect
to $\log S_t$ is
$$
  \frac{d\log E_t}{d\log S_t}
  = \underbrace{\frac{d\pi_t}{d\log S_t}}_{\gamma_{\pi S}}
  + \underbrace{\frac{A_t}{A_t - D_t}}_{\mathrm{ILE}_t}.
$$
\end{proof}

\paragraph{Economic intuition.}
When BTC rises, the premium typically \emph{compresses} ($\gamma_{\pi S} < 0$):
the market discounts the option as the stock already moved.  This dampening
is analogous to the delta of an option declining as the underlying moves away
from the strike.  With calibrated values,
$$
  \mathrm{TEE}_0 = \gamma_{\pi S} + \mathrm{ILE}_0
  = -3.60 + 1.49 = -2.11,
$$
indicating that net equity actually moves \emph{inversely} to BTC at the margin---
a regime where premium compression dominates leverage.

\begin{remark}
The negative TEE is consistent with the empirical observation that in periods of
rapidly falling premium, MSTR equity can underperform BTC even on up-days.
The sign of $\gamma_{\pi S}$ is regime-dependent; in premium-expansion regimes
$\gamma_{\pi S} > 0$ and TEE exceeds ILE.
\end{remark}

%% ------------------------------------------------------------------
\subsection{PMRI as Theta}\label{ssec:pmri-theta}

\begin{proposition}[PMRI $\equiv$ Premium Theta]\label{prop:pmri-theta}
Under the OU dynamics $d\pi_t = \kappa(\theta - \pi_t)\,dt + \sigma_\pi\,dW_t^{\pi}$,
the standardised expected premium decay is
$$
  \Theta_{\pi}(t)
  = \frac{\pi_t - \theta}{\sigma_{\mathrm{stat}}}
  = \mathrm{PMRI}_t,
  \qquad
  \sigma_{\mathrm{stat}} = \frac{\sigma_\pi}{\sqrt{2\kappa}}.
$$
\end{proposition}

\begin{proof}
The stationary distribution of the OU process is $\pi_\infty \sim
\mathcal{N}\!\bigl(\theta,\,\sigma_\pi^2/(2\kappa)\bigr)$.
The z-score of the current premium relative to this distribution is
$(\pi_0 - \theta)/\sigma_{\mathrm{stat}}$, which is the PMRI by definition.
The expected instantaneous drift is $\kappa(\theta - \pi_t)$; normalising by
the stationary standard deviation yields a unit-free measure of the ``carry
cost'' of the current premium level.
\end{proof}

\paragraph{Economic intuition.}
Theta in options measures time decay---the cost of holding a position as time
passes.  PMRI plays the same role for the accumulation option: a high
premium ($\mathrm{PMRI} \gg 0$) signals that the premium is likely to decay,
eroding the option value.  A negative PMRI indicates the premium is cheap
relative to its long-run level, conferring positive carry.

With calibrated parameters $\kappa = 5.21$, $\theta = 1.08$,
$\sigma_\pi = 3.79$, and $\pi_0 = 0.24$:
$$
  \sigma_{\mathrm{stat}} = \frac{3.79}{\sqrt{2 \times 5.21}} = 1.175,
  \qquad
  \mathrm{PMRI}_0 = \frac{0.24 - 1.08}{1.175} = -0.71.
$$
The negative value implies the premium is \emph{below} its long-run mean,
conferring positive expected carry---the accumulation option is ``cheap.''

%% ------------------------------------------------------------------
\subsection{IBGR as Moneyness}\label{ssec:ibgr-moneyness}

\begin{proposition}[IBGR $\equiv$ Moneyness]\label{prop:ibgr-moneyness}
The implied BTC growth rate
$$
  \mathrm{IBGR}_t = \mu_H(t)
  = \frac{1}{H_t}\bigl(\alpha\,\pi_t^+ + \lambda_M\,\mathbb{E}[\Delta H]\bigr)
$$
measures the \emph{moneyness} of the accumulation option: how deeply in the
money the accretive strategy is.
\end{proposition}

\begin{proof}
The accumulation option is valuable precisely when management can grow BTC
holdings at a rate exceeding what the market has priced.  IBGR quantifies
this rate.  When $\mathrm{IBGR} > 0$, the option is in the money---each
financing event is expected to add BTC.  A higher IBGR corresponds to
deeper moneyness, as the expected BTC accretion per unit time is larger,
making the option more valuable.

Formally, the value of the accumulation option scales with IBGR through the
premium dynamics: $V_{\mathrm{acc}} \propto \mathbb{E}\bigl[\int_0^T
e^{-rs}(e^{\pi_s}-1)(A_s - D_s)^+\,ds\bigr]$, and $\pi_s$ is sustained at
high levels precisely when holdings grow fast (the flywheel), i.e.\ when
$\mu_H$ is large.
\end{proof}

\paragraph{Economic intuition.}
Just as an option's moneyness (ratio of underlying to strike) determines
its intrinsic value, IBGR determines the ``intrinsic accretion'' embedded in
the stock.  With the current calibration:
$$
  \mathrm{IBGR}_0
  = \frac{1}{761{,}068}\bigl(0 + 19.05 \times 7{,}044\bigr)
  = 17.6\%\;\text{annualised},
$$
indicating the accumulation option is significantly in the money.

%% ------------------------------------------------------------------
\subsection{IFRD as Exercise Cost}\label{ssec:ifrd-strike}

\begin{proposition}[IFRD $\equiv$ Exercise Cost]\label{prop:ifrd-strike}
The implied funding requirement distribution
$$
  G_T = \int_0^T S_u\,dH_u
$$
is the \emph{exercise cost} of the accumulation option---the total capital
that must be deployed to realise the implied BTC growth.
\end{proposition}

\begin{proof}
In a standard call option, the strike price $K$ is the cost of exercising.
For the accumulation option, there is no single exercise event; instead,
management continuously ``exercises'' by purchasing BTC.  The cumulative
dollar cost $G_T$ is the stochastic analogue of the strike.  Its distribution
under $\mathbb{P}$ (the IFRD) characterises the uncertainty in this
exercise cost.

From the simulation under calibrated dynamics:
\begin{align*}
  \mathbb{E}[G_{3\mathrm{y}}] &= \$35.3\text{B}, \\
  G_{3\mathrm{y}}^{95\%} &= \$71.8\text{B}.
\end{align*}
\end{proof}

\paragraph{Economic intuition.}
The IFRD answers: ``How much capital must be raised to keep the flywheel
spinning?''  Like a strike price, it determines whether exercising the option
is economically rational.  If the premium-adjusted equity issuance can cover
$G_T$ without excessive dilution, the option is worth exercising.  The
95th-percentile IFRD of \$71.8B flags the tail scenario where capital markets
must remain extremely accommodative.

%% ------------------------------------------------------------------
\subsection{Summary of the Greek Mapping}

\begin{table}[ht]
\centering
\caption{Structural indicators as option Greeks.}
\label{tab:greek-map}
\begin{tabular}{@{}llll@{}}
\toprule
\textbf{Indicator} & \textbf{Greek} & \textbf{Formula} & \textbf{Calibrated} \\
\midrule
ILE $= A/(A-D)$ & Delta ($\Delta_{\mathrm{lev}}$)
  & $\partial\log E / \partial\log S\big|_\pi$
  & $1.49$ \\[4pt]
TEE $= \gamma_{\pi S} + \text{ILE}$ & Adjusted Delta ($\Delta_{\mathrm{tot}}$)
  & $d\log E / d\log S$
  & $-2.11$ \\[4pt]
PMRI $= (\pi-\theta)/\sigma_{\mathrm{stat}}$ & Theta ($\Theta_\pi$)
  & Premium carry cost
  & $-0.71$ \\[4pt]
IBGR $= \mu_H$ & Moneyness
  & $(\alpha\pi^+ + \lambda_M\mathbb{E}[\Delta H])/H$
  & $17.6\%$ \\[4pt]
IFRD $= G_T$ & Exercise cost
  & $\int S\,dH$
  & \$35.3B (mean) \\
\bottomrule
\end{tabular}
\end{table}

%% ------------------------------------------------------------------
\subsection{New Greeks: Vega and Rho}\label{ssec:new-greeks}

The standard option-Greek framework suggests two additional sensitivities
absent from the original indicator set.

\subsubsection{Vega: Sensitivity to BTC Volatility}

\begin{definition}[Premium Vega]\label{def:vega}
Define
$$
  \mathcal{V}_\pi
  = \frac{\partial\,\pi^*}{\partial\,\sigma_S},
$$
where $\pi^* = \theta$ is the long-run fair premium implied by the OU
equilibrium.
\end{definition}

\begin{proposition}[Positive Vega]\label{prop:vega}
The accumulation option has positive vega: higher BTC volatility increases
the fair premium.
\end{proposition}

\begin{proof}
The accumulation option's value derives from management's ability to exploit
periods of high premium.  Consider the premium dynamics
$d\pi_t = \kappa(\theta - \pi_t)\,dt + \sigma_\pi\,dW_t^\pi$ with
$\mathrm{corr}(dW^S, dW^\pi) = \rho$.

Higher BTC volatility $\sigma_S$ affects the system through two channels:
\begin{enumerate}
  \item \textbf{Option convexity.}  The equity payoff $(A_t - D_t)^+$ is
    convex in $S_t$.  By Jensen's inequality, higher $\sigma_S$ increases
    the expected value of this payoff, justifying a higher equilibrium
    premium.  Formally, for the leveraged claim:
    $$
      \mathbb{E}\bigl[(H S_T - D)^+\bigr]
      \;\text{is increasing in } \sigma_S,
    $$
    by standard option-pricing monotonicity.

  \item \textbf{Flywheel optionality.}  Higher $\sigma_S$ generates more
    frequent episodes of extreme BTC appreciation, during which the premium
    expands ($\pi_t \uparrow$) and management issues equity to buy BTC.
    The drift term $\alpha\,\pi_t^+$ in holdings dynamics implies that
    higher vol creates more accumulation opportunities.
\end{enumerate}

Both channels push the fair premium $\pi^*$ upward with $\sigma_S$, yielding
$\mathcal{V}_\pi > 0$.
\end{proof}

\paragraph{Calibration.}
To estimate $\mathcal{V}_\pi$ numerically, we perturb $\sigma_S$ by
$\pm 5\%$ around the calibrated value $\hat\sigma_S = 0.487$ and re-run the
Monte Carlo simulation to obtain the equilibrium premium in each scenario.
The finite-difference estimate is
$$
  \hat{\mathcal{V}}_\pi
  = \frac{\pi^*(\sigma_S + \delta) - \pi^*(\sigma_S - \delta)}{2\delta}.
$$

\paragraph{Interpretation.}
Positive vega explains why MSTR's premium tends to be elevated in
high-volatility BTC regimes: the accumulation option becomes more
valuable when the underlying is more volatile, exactly as for standard
options.

\subsubsection{Rho: Sensitivity to Interest Rates}

\begin{definition}[Premium Rho]\label{def:rho-premium}
Define
$$
  \mathcal{R}_\pi
  = \frac{\partial\,\pi^*}{\partial\,r},
$$
where $r$ is the risk-free rate (or, more precisely, MSTR's cost of debt
capital).
\end{definition}

\begin{proposition}[Negative Rho]\label{prop:rho-negative}
The accumulation option has negative rho: higher interest rates decrease
the fair premium.
\end{proposition}

\begin{proof}
Higher rates affect the accumulation option through three channels:
\begin{enumerate}
  \item \textbf{Debt servicing cost.}  MSTR's debt $D_t$ carries interest.
    Higher $r$ increases the carrying cost of leverage, reducing the net
    benefit of the leveraged BTC position and narrowing the region where
    $A_t > D_t$.

  \item \textbf{Financing cost.}  The flywheel requires capital market
    access.  Higher rates increase the cost of new debt issuance and
    reduce the attractiveness of convertible offerings, slowing the
    expected BTC accumulation rate $\mu_H$.

  \item \textbf{Discount rate.}  The present value of future accumulation
    benefits is discounted at a higher rate, reducing $V_{\mathrm{acc}}$
    and hence the equilibrium premium.
\end{enumerate}

Formally, in the reduced-form OU model, $\theta$ is determined in
equilibrium by market pricing of these factors.  Holding other parameters
fixed, an increase in $r$ compresses the equilibrium premium:
$\partial\theta/\partial r < 0$.  Since $\pi^* = \theta$, we obtain
$\mathcal{R}_\pi < 0$.
\end{proof}

\paragraph{Calibration.}
As with vega, $\mathcal{R}_\pi$ can be estimated by perturbing the
risk-free rate in the simulation.  The preferred dividend obligations
(\$977M annually at current calibration) make the rate sensitivity
particularly acute: higher rates increase the probability that the
dividend coverage ratio falls below one, threatening the sustainability
of the capital structure.

\begin{remark}
The sign conventions match standard equity-option Greeks: a leveraged
long position has positive delta, positive vega, and (for puts or
leveraged structures with debt) negative rho.  The accumulation option
adds the novel feature that theta is endogenous---management can
influence it through the flywheel.
\end{remark}
